We present \texttt{micrograd}, an open-source FEniCSx framework for
density-based topology optimisation of microfluidic devices governed by
coupled Brinkman--convection-diffusion transport.
The principal computational challenges are (i)~efficient sensitivity
computation for a multi-objective cost functional that simultaneously
penalises viscous dissipation and outlet concentration mismatch, and
(ii)~robust stabilisation of the convection-dominated transport equation
and its adjoint at high cell P\'eclet numbers.

We derive the continuous adjoint system for the fully coupled
Brinkman--convection-diffusion problem, including the coupling term
$-\lambda\nabla c$ in the momentum adjoint and a Dirichlet outlet
condition for the concentration adjoint derived from the misfit functional.
The implementation in FEniCSx~(DOLFINx~v0.7.3) uses Taylor--Hood
P2/P1 elements for flow and P1+SUPG elements for transport;
the SUPG parameter is evaluated element-wise from the local P\'eclet
number and applied consistently to both the forward and adjoint transport
equations.
Helmholtz PDE filtering, Heaviside projection with
$\beta$-continuation to $\beta=\betaMax$, and logarithmic
$\alpha_{\max}$ ramping to $\alphaMaxFinal$ together drive the
design toward low gray-zone fractions.

Functional smoothness is confirmed by a finite-difference Taylor remainder test
(slope $\approx 2.0$ on both coarse and primary meshes).
The continuously assembled adjoint captures the correct descent direction
(Pearson correlation $\geq 0.80$ with finite-difference gradients);
a discrete adjoint with exact gradient magnitude is identified as future work.
Manufactured-solution tests confirm the expected spatial convergence
rates for the forward and adjoint PDE solvers.

As a proof-of-concept on an $80\times20$ finite-element mesh, the framework is applied to the inverse design of
a two-inlet microfluidic concentration gradient generator.
The optimised design achieves an outlet RMSE of $\rmseTopOpt$ for a
linear target profile and reduces hydraulic resistance by $\reduction\%$
relative to the classic Christmas-tree network~\cite{jeon2000}.
The complete pipeline runs on a standard laptop in under thirty minutes
and is fully reproducible via Docker and continuous integration.
